% Teil II: durchgerechneter Weg, Teile a--b.

\subsection*{Gegeben und gesucht}
$U(r)=-GMm/r$ mit $r=\sqrt{x^2+y^2+z^2}$. Gesucht: (a) $\vec F=-\nabla U$
durch explizite Gradientenbildung, (b)(Bonus) $\nabla\times\vec F$ und
Aussage zur Konservativitaet.

\subsection*{Ansatz}
Partielle Ableitungen von $U(r)$ über die Kettenregel mit
$\partial r/\partial x=x/r$ etc.; sodann der Gradientenvektor. Für die
Rotation nutzt man die Identitaet $\nabla\times(-\nabla U)=\vec 0$.

\subsection*{Rechnung}
\textbf{a) Gradient und Kraft.} Mit
$\mathrm dU/\mathrm dr=GMm/r^2$ (da $\mathrm d(-1/r)/\mathrm dr=1/r^2$)
gilt
\begin{equation}
  \frac{\partial U}{\partial x}=\frac{GMm}{r^2}\cdot\frac{x}{r}
  =\frac{GMm\,x}{r^3},
\end{equation}
analog für $y,z$. Damit
\begin{equation}
  \nabla U=\frac{GMm}{r^3}(x,y,z)=\frac{GMm}{r^2}\,\hat e_r,
\end{equation}
und mit dem Vorzeichen
\begin{equation}
  \boxed{\;\vec F=-\nabla U=-\frac{GMm}{r^3}(x,y,z)
  =-\frac{GMm}{r^2}\,\hat e_r\;}.
\end{equation}
Die Kraft zeigt radial nach innen (Anziehung), Betrag $GMm/r^2$.

\textbf{b) Rotation (Bonus).} Da $\vec F=-\nabla U$ gilt, ist
\begin{equation}
  \nabla\times\vec F=-\nabla\times(\nabla U)=\vec 0
  \quad(\text{per Komponente prüfbar})\;.
\end{equation}
Das Gravitationsfeld ist also wirbelfrei. Mit dem Satz von Stokes folgt
damit für jeden geschlossenen Weg $\oint\vec F\cdot\mathrm d\vec r=0$,
d. h. das Feld ist \emph{konservativ}.

\begin{ergebnis}
$\vec F=-(GMm/r^3)(x,y,z)=-(GMm/r^2)\hat e_r$; $\nabla\times\vec F=\vec 0$,
das Gravitationsfeld ist konservativ, das Wegintegral ist wegunabhängig.
\end{ergebnis}

\begin{kontrolle}
\textbf{Querprobe (Radialsymmetrie):} Für $\vec F=-GMm\,\hat e_r/r^2$ sind
alle Ableitungen entlang der Komponenten konsistent (z. B. $\partial_y F_z$-
Terme heben sich mit $\partial_z F_y$-Termen), also $\nabla\times\vec F=\vec 0$.
\checkmark

\textbf{Vorzeichen:} $F<0$ in $\hat e_r$-Richtung bedeutet Attraktion,
wie für die Gravitation erwartet. \checkmark

\textbf{Einheiten:} $\nabla U$: [J/m] = [N], was zu einer Kraft passt. \checkmark
\end{kontrolle}

\begin{merksatz}
Gradientenfelder $\vec F=-\nabla U$ sind automatisch rotationfrei
($\nabla\times\vec F=\vec 0$) und damit konservativ: Die Arbeit längs
eines Weges hängt nur von den Endpunkten ab. Die Gravitation zweigt als
radialsymmetrisches Potential $(GMm/r)$ genau dieses Bild.
\end{merksatz}
